% Fontes do capítulo: uma entrada por chave \src{…}. Ordem = ordem de primeira citação (aproximada).
% Fonte primária ao lado da secundária frágil sempre que localizada (A5, E-R3-5).
% Formato: Autor, \emph{Obra}, local (trad./ed.). \url{…}  ·  artigos com DOI.
\chapter*{Fontes}
\addcontentsline{toc}{chapter}{Fontes}
\markboth{Fontes}{Fontes}
\small
Os números sobrescritos no texto remetem a esta lista. As questões trazem a procedência no próprio cabeçalho. Os endereços da web foram acessados em novembro de 2026, e em todas as páginas da web (Wikipédia, enciclopédias e repositórios) vale a versão dessa data. Quando a fonte acessível é secundária, a obra original vem citada ao lado, sempre que localizada.

\begingroup
\makeatletter
% a lista entra logo abaixo do título "Fontes": sem novo capítulo nem cabeçalho "Referências"
\renewcommand{\chapter}{\@ifstar{\@gobble}{\@gobble}}
\renewcommand{\@mkboth}[2]{}
\makeatother
\begin{thebibliography}{999}\setlength{\itemsep}{0.2em}
\bibitem{fme} G.~Iezzi et al., \emph{Fundamentos de Matemática Elementar}, vol.~1, cap.~II, \enquote{Conjuntos}. São Paulo: Atual.
\end{thebibliography}
\endgroup
